% Original: PDF strana 25, donji slajd.
\slajdid{02-s050}
\begin{frame}[label=02-s050]{Lažne nule i jedinice — glič}
\setlength{\parskip}{0pt}
% element: 02-s050:e01
Kašnjenja logičkih kola razlikuju se i zavise od fabrikacije, temperature, starosti komponenti i drugih uslova. Zato se pri promeni ulaznih signala mogu pojaviti neželjeni efekti na izlazu.
\par\vspace{2mm}
\begin{columns}[T,onlytextwidth]
\begin{column}{.48\textwidth}
\Oznaka{Lažna nula}
Izlaz je 1 i posle promene ulaza treba da ostane 1. Zbog prostiranja signala ipak se pojavi \alert{\textbf{kratkotrajna nula}}.
\end{column}
\begin{column}{.48\textwidth}
\Oznaka{Lažna jedinica}
Izlaz je 0 i posle promene ulaza treba da ostane 0. Zbog prostiranja signala ipak se pojavi \alert{\textbf{kratkotrajna jedinica}}.
\end{column}
\end{columns}
\par\vspace{2mm}
Ovi kratkotrajni neželjeni impulsi nazivaju se \alert{\textbf{glič}} (engl. \emph{glitch}).
\par\vspace{1mm}
% element: 02-s050:e02
\Oznaka{Posmatrajmo primer kombinacione mreže: \(F=C\bar B+BA\).}
{\centering\begin{tikzpicture}[circuit logic US,DEoznaka,x=1cm,y=1cm,
 gate/.style={draw,logic gate inputs=nn,minimum width=12mm,minimum height=9mm,scale=.6}]
\node[and gate US,gate] (u) at (2,1.1) {};
\node[and gate US,gate] (d) at (2,-.1) {};
\node[or gate US,gate] (f) at (3.8,.5) {};
\node[not gate US,draw,minimum width=8mm,minimum height=8mm,scale=.6] (i) at (.6,.55) {};
\draw[DEvod] (u.input 1)--(-.5,0 |- u.input 1) node[left] {$C$};
\draw[DEvod] (i.input)--(-.5,.55) node[left] {$B$};
\draw[DEvod] (i.output)--(1.25,.55)|-(u.input 2);
\draw[DEvod] (-.2,.55)|-(d.input 1);
\node[junction] at (-.2,.55) {};
\draw[DEvod] (d.input 2)--(-.5,0 |- d.input 2) node[left] {$A$};
\draw[DEvod] (u.output)--(2.8,1.1)|-(f.input 1);
\draw[DEvod] (d.output)--(2.8,-.1)|-(f.input 2);
\draw[DEvod] (f.output)--++(.4,0) node[right] {$F$};
\end{tikzpicture}
\par}
\end{frame}
